\chapter{Marco Referencial}
\section{Marco Teórico}
La creciente adopción de la gestión por procesos ha estado estrechamente vinculada con la transformación digital, fenómeno que ha modificado la manera en que las organizaciones diseñan, ejecutan y controlan sus actividades. En este contexto, los procesos empresariales se consolidan como el eje central de la estrategia organizacional, ya que permiten integrar tecnologías digitales orientadas a optimizar el desempeño institucional, incrementar la productividad y fortalecer la competitividad. Al respecto, Gabryelczyk, Sipior y Biernikowicz (2022) sostienen que la transformación digital ha impulsado la evolución del Business Process Management (BPM), proporcionando mecanismos que incrementan la agilidad operativa, mejoran la coordinación entre las diferentes áreas funcionales y permiten responder con mayor rapidez a los cambios del entorno empresarial.

Desde esta perspectiva, la transformación digital no debe entenderse únicamente como la incorporación de nuevas herramientas tecnológicas, sino como un proceso integral de cambio organizacional que modifica la forma en que las empresas generan valor. Carmona Arbañil (2023) afirma que la implementación de tecnologías digitales favorece la optimización de las actividades organizacionales mediante la automatización de procesos, el fortalecimiento de la productividad y la mejora continua de la calidad de los servicios. En consecuencia, la integración entre la gestión por procesos y las tecnologías digitales constituye un elemento estratégico que impulsa la innovación y facilita el desarrollo sostenible de las organizaciones en entornos altamente competitivos.

Como resultado de esta evolución tecnológica, la automatización de procesos empresariales se ha convertido en una herramienta indispensable para incrementar la eficiencia operativa y garantizar la estandarización de las actividades organizacionales. La automatización permite ejecutar tareas repetitivas con una mínima intervención humana, reduciendo la variabilidad de los procesos, optimizando el uso de los recursos y disminuyendo la probabilidad de errores operativos. En este sentido, Tinoco-León (2017) señala que la automatización de procesos de negocio mejora la capacidad de adaptación de las organizaciones, ya que reduce los tiempos de respuesta frente a las demandas del entorno y favorece el desarrollo de procedimientos estandarizados que incrementan la eficiencia de las operaciones.

En concordancia con lo anterior, Cervera Martínez (2025) sostiene que la automatización de procesos administrativos mediante sistemas digitales estructurados permite establecer flujos de trabajo más organizados, controlados y eficientes. Asimismo, destaca que estos sistemas incorporan mecanismos de seguimiento y supervisión que fortalecen el control interno, garantizan la precisión de las actividades ejecutadas y mejoran la administración de la información organizacional. De esta manera, la automatización no solo representa una alternativa para reducir costos operativos, sino que también constituye un factor clave para fortalecer la calidad de los procesos y optimizar la gestión empresarial.

La evolución de las tecnologías digitales ha favorecido, además, la incorporación de la Inteligencia Artificial (IA) como una herramienta complementaria a la automatización tradicional. A diferencia de los sistemas convencionales, la IA posee la capacidad de analizar grandes volúmenes de datos, identificar patrones de comportamiento, aprender de la información procesada y apoyar la toma de decisiones en tiempo real. Estas capacidades permiten desarrollar procesos más inteligentes, flexibles y adaptativos, incrementando la eficiencia organizacional y fortaleciendo la capacidad de respuesta frente a escenarios dinámicos y altamente competitivos.

Los beneficios derivados de la integración entre automatización e Inteligencia Artificial han sido evidenciados en diversos estudios recientes. Navas Espin, Zambrano Silva, Pincay Bohórquez, Lizarzaburu Mora y Segura Torres (2026) demostraron que la implementación conjunta de estas tecnologías genera mejoras significativas en los indicadores de eficiencia operacional. Entre los principales resultados obtenidos destacan una reducción del 76\,\% en las tasas de error, un incremento del 50\,\% en la productividad por operario y una disminución del 60\,\% en el tiempo del ciclo de producción. Estos hallazgos evidencian que la incorporación de tecnologías inteligentes no solo optimiza el desempeño operativo de las organizaciones, sino que también fortalece su competitividad, facilita la toma de decisiones basada en datos y contribuye al desarrollo de procesos más eficientes, confiables y orientados hacia la mejora continua.
La transformación digital ha impulsado una nueva concepción de la gestión organizacional, en la que los procesos empresariales constituyen el eje central para alcanzar mayores niveles de eficiencia, productividad y competitividad. En este contexto, las organizaciones han incorporado tecnologías digitales que les permiten optimizar sus operaciones, fortalecer la coordinación entre las diferentes áreas funcionales y responder con mayor rapidez a los cambios del entorno. Gabryelczyk, Sipior y Biernikowicz (2022) sostienen que la adopción del Business Process Management (BPM) se ha fortalecido como consecuencia de la transformación digital, ya que proporciona mecanismos que incrementan la agilidad operativa, mejoran la integración de los procesos y favorecen la capacidad de adaptación frente a escenarios empresariales dinámicos.

Desde esta perspectiva, la transformación digital no se limita a la incorporación de herramientas tecnológicas, sino que implica una transformación integral de la organización mediante la redefinición de sus procesos y modelos de gestión. En este sentido, Carmona Arbañil (2023) afirma que la implementación de tecnologías digitales contribuye a optimizar las actividades organizacionales, incrementar la productividad, mejorar la calidad de los servicios e impulsar la innovación empresarial. En consecuencia, la integración entre la gestión por procesos y las tecnologías digitales constituye un elemento estratégico para fortalecer la competitividad organizacional y garantizar la mejora continua.

Como resultado de esta evolución tecnológica, la automatización de procesos empresariales ha adquirido una importancia significativa dentro de las organizaciones modernas. La automatización permite ejecutar tareas repetitivas mediante sistemas tecnológicos que reducen la intervención humana, disminuyen la variabilidad de los procesos y optimizan el aprovechamiento de los recursos disponibles. De acuerdo con Tinoco-León (2017), la automatización de procesos de negocio incrementa la capacidad de adaptación de las organizaciones, ya que permite responder con mayor rapidez a las exigencias del entorno, reducir errores operativos y estandarizar los procedimientos administrativos y productivos.

En complemento, Cervera Martínez (2025) sostiene que la automatización de procesos administrativos mediante sistemas digitales estructurados favorece el desarrollo de flujos de trabajo más organizados y controlados, incorporando mecanismos de supervisión que garantizan la precisión de las actividades realizadas. Esta integración tecnológica fortalece el control interno, mejora la gestión de la información y contribuye a optimizar la utilización de los recursos organizacionales, generando beneficios tanto en la eficiencia operativa como en la calidad de los servicios.

La evolución de la automatización ha propiciado la incorporación de la Inteligencia Artificial (IA) dentro de los procesos empresariales, permitiendo que los sistemas tecnológicos no solo ejecuten tareas previamente programadas, sino que también aprendan a partir de los datos, identifiquen patrones y apoyen la toma de decisiones en tiempo real. Estas capacidades convierten a la IA en un recurso estratégico para desarrollar procesos más inteligentes, flexibles y adaptativos, fortaleciendo la capacidad de respuesta de las organizaciones frente a entornos cada vez más cambiantes.

Los beneficios derivados de la integración entre automatización e Inteligencia Artificial han sido evidenciados por Navas Espin, Zambrano Silva, Pincay Bohórquez, Lizarzaburu Mora y Segura Torres (2026), quienes demostraron que la implementación de estas tecnologías produce mejoras significativas en los indicadores de eficiencia operacional. Entre los principales resultados obtenidos destacan una reducción del 76\,\% en las tasas de error, un incremento del 50\,\% en la productividad por operario y una disminución del 60\,\% en el tiempo del ciclo de producción. Estos resultados evidencian que la incorporación de tecnologías inteligentes representa un factor determinante para fortalecer la competitividad empresarial, optimizar los procesos y consolidar una cultura organizacional orientada hacia la mejora continua.
La evolución de los sistemas de planificación de recursos empresariales (\textit{Enterprise Resource Planning}, ERP) ha permitido que las organizaciones sustituyan modelos tradicionales de gestión por plataformas integradas capaces de centralizar la información y coordinar los diferentes procesos organizacionales. Estos sistemas facilitan la administración de los recursos empresariales mediante la integración de áreas como producción, inventarios, compras, ventas, finanzas y recursos humanos, proporcionando información oportuna para apoyar la toma de decisiones. En este sentido, Sánchez-Calderón et al. (2024) señalan que la evolución de los sistemas ERP responde a las nuevas exigencias del entorno digital, permitiendo mejorar la integración de la información, optimizar la gestión de los recursos y fortalecer la eficiencia de las operaciones empresariales.

De manera complementaria, Algaba Berro et al. (2017) sostienen que la implementación de plataformas ERP constituye una estrategia organizacional que favorece la administración centralizada de la información y mejora la competitividad mediante la integración de las diferentes áreas funcionales de la empresa. Los autores destacan que estos sistemas reducen la duplicidad de información, optimizan los procesos administrativos y fortalecen el control de las operaciones, convirtiéndose en un soporte fundamental para la gestión empresarial.

En los últimos años, el software empresarial de código abierto ha adquirido una importancia creciente debido a su capacidad para proporcionar soluciones tecnológicas flexibles, escalables y de menor costo para las organizaciones, especialmente para las pequeñas y medianas empresas. Bhattacharyya y Dan (2014) sostienen que los sistemas ERP de código abierto representan una alternativa viable para integrar procesos organizacionales sin requerir inversiones elevadas, permitiendo adaptar el software a las necesidades específicas de cada empresa y facilitando la optimización de los recursos tecnológicos.

De forma complementaria, Quimis Sánchez y Sánchez Rodríguez (2017) manifiestan que las aplicaciones informáticas especializadas permiten reemplazar procedimientos manuales por procesos digitales más eficientes, mejorando la precisión en la administración de la información y reduciendo los tiempos de ejecución de las actividades organizacionales. En conjunto, estos autores coinciden en que las soluciones tecnológicas basadas en software empresarial fortalecen la eficiencia administrativa y contribuyen al desarrollo de organizaciones más competitivas y orientadas hacia la transformación digital.

La evaluación permanente del desempeño organizacional requiere indicadores que permitan medir objetivamente la eficiencia de los procesos. Entre ellos, la Eficiencia General de los Equipos (\textit{Overall Equipment Effectiveness}, OEE) constituye uno de los indicadores más utilizados para evaluar el nivel de aprovechamiento de los recursos productivos mediante la integración de tres dimensiones fundamentales: disponibilidad, rendimiento y calidad. El análisis conjunto de estos componentes facilita la identificación de pérdidas operacionales, permitiendo establecer estrategias orientadas a incrementar la productividad y fortalecer los procesos de mejora continua.

Para representar, documentar y analizar los procesos organizacionales, la Notación para el Modelado de Procesos de Negocio (\textit{Business Process Model and Notation}, BPMN) se ha consolidado como el estándar internacional más utilizado. Esta metodología emplea una representación gráfica basada en eventos, actividades, compuertas y flujos de secuencia que facilita la comprensión de los procesos tanto para los responsables del negocio como para los desarrolladores de sistemas. Su aplicación permite identificar oportunidades de mejora, optimizar el diseño de los procesos y facilitar la implementación de soluciones tecnológicas orientadas a la automatización.

Como complemento al análisis de procesos, la Inteligencia de Negocios (\textit{Business Intelligence}, BI) proporciona herramientas que permiten recopilar, integrar y analizar grandes volúmenes de datos provenientes de las diferentes áreas de la organización. Entre sus principales componentes se encuentran los \textit{data warehouses}, los \textit{data marts} y los tableros de control (\textit{dashboards}), los cuales facilitan la visualización de indicadores clave de desempeño (KPI) y apoyan la toma de decisiones basada en evidencia. La integración entre sistemas ERP, metodologías de modelado de procesos, indicadores de desempeño y herramientas de Inteligencia de Negocios constituye un soporte tecnológico fundamental para fortalecer la gestión organizacional, incrementar la eficiencia operativa y consolidar procesos de mejora continua orientados hacia la competitividad empresarial.

\section{Marco Conceptual}

\textbf{PYMES (Pequeñas y Medianas Empresas):} Son unidades económicas productivas de tamaño reducido o mediano que desempeñan actividades comerciales, industriales o de servicios. Se caracterizan por contar con un número limitado de empleados y recursos financieros, lo que influye en su capacidad de inversión tecnológica y en la adopción de sistemas de información avanzados.

\textbf{BPMM (Business Process Management Maturity / Gestión de la Madurez de Procesos de Negocio):} Enfoque que evalúa el nivel de desarrollo y capacidad de una organización para gestionar, controlar y mejorar sus procesos de negocio de manera sistemática. Permite identificar el grado de madurez de los procesos empresariales y establecer estrategias de mejora continua orientadas a la optimización operativa y la transformación digital.


\textbf{Aplicación Web:} Software desarrollado para ejecutarse a través de un navegador web mediante protocolos de Internet. Permite a los usuarios acceder a funcionalidades empresariales sin necesidad de instalar programas adicionales en sus dispositivos.

\textbf{Base de Datos:} Conjunto organizado de información almacenada de manera estructurada para facilitar su consulta, actualización y administración.

\textbf{Sistema de Información:} Conjunto de elementos tecnológicos, procedimientos y recursos humanos que permiten recopilar, procesar, almacenar y distribuir información para apoyar la toma de decisiones.

\textbf{Automatización de Procesos:} Uso de tecnologías informáticas para ejecutar actividades de manera automática, reduciendo la intervención manual y mejorando la eficiencia operativa.

\textbf{ERP (Enterprise Resource Planning):} Sistema de planificación de recursos empresariales que integra diferentes áreas de una organización, como ventas, inventario, producción, finanzas y recursos humanos.

\textbf{BPMS (Business Process Management System):} Plataforma tecnológica utilizada para modelar, automatizar, ejecutar, monitorear y optimizar procesos de negocio.

\textbf{Transformación Digital:} Proceso mediante el cual una organización incorpora tecnologías digitales para mejorar sus operaciones, generar valor y adaptarse a nuevas exigencias del mercado.

\textbf{Proforma:} Documento comercial utilizado para presentar una cotización preliminar de productos o servicios antes de concretar una venta.

\textbf{Liquidación:} Documento que registra el cálculo final de valores económicos asociados a una transacción comercial.

\textbf{API (Application Programming Interface):} Conjunto de reglas y mecanismos que permiten la comunicación entre diferentes aplicaciones o sistemas.

\textbf{Framework:} Estructura de desarrollo de software que proporciona herramientas y componentes reutilizables para construir aplicaciones de forma eficiente.

\textbf{ORM (Object Relational Mapping):} Técnica que permite interactuar con bases de datos mediante objetos de programación sin escribir consultas SQL complejas.

\textbf{Escalabilidad:} Capacidad de un sistema para soportar un aumento en la cantidad de usuarios, datos o procesos sin afectar significativamente su rendimiento.

\section{Marco Contextual}
La fábrica Confort Muebles inició sus actividades en el año 2009 en la ciudad de
Cuenca, Ecuador, dedicándose a la fabricación y comercialización de muebles para el
hogar y oficina, así como a la elaboración de muebles a medida según los requerimientos
de sus clientes. Desde su creación, la empresa ha buscado ofrecer productos de calidad,
adaptándose a las necesidades del mercado local.
La empresa es de propiedad de María Lucero, mientras que la administración se
encuentra a cargo de Luis Luzuriaga, quien es responsable de coordinar las actividades
administrativas y operativas. Entre sus funciones se encuentran la supervisión de los
procesos internos, el abastecimiento de materiales necesarios para la producción, la
revisión de las proformas y el control de las actividades comerciales de la empresa.
En el área de producción, Confort Muebles cuenta con un equipo conformado
por seis trabajadores, quienes se encargan de la fabricación de los diferentes productos.
El proceso productivo depende de la disponibilidad de materiales, por lo que el adminis
trador realiza un seguimiento permanente del inventario para garantizar la continuidad


Actualmente, la gestión de la información relacionada con las proformas, clien
tes, inventario y liquidaciones se realiza principalmente de forma manual, utilizando
documentos físicos y archivos independientes. Esta modalidad de trabajo ha ocasionado
inconvenientes como errores en el registro de datos, duplicidad de documentos, pérdida
de información y dificultades para localizar los registros cuando son requeridos.
Asimismo, la empresa no dispone de un sistema informático que permita centra
lizar la información y automatizar los procesos administrativos. Como consecuencia, el
seguimiento de las cotizaciones y de los clientes resulta más complejo, incrementando
los tiempos de atención y afectando la eficiencia de las actividades comerciales.

Ante esta situación, surge la necesidad de implementar una aplicación web
que permita automatizar la gestión de proformas, liquidaciones, clientes, productos
e inventario. La incorporación de esta solución tecnológica permitirá centralizar la
información en una única base de datos, reducir errores derivados de los procesos
manuales, agilizar la elaboración y seguimiento de las cotizaciones, mejorar el control
administrativo y fortalecer la toma de decisiones basada en información confiable y
actualizada.
La implementación de este sistema contribuirá a optimizar los procesos admi
nistrativos y comerciales de Confort Muebles, ofreciendo una atención más rápida y
eficiente a los clientes, mejorando la organización de la información y fortaleciendo la
competitividad de la fabrica dentro del sector manufacturero de la ciudad de Cuenca